\documentclass[11pt]{amsart}
\usepackage{amsmath,amssymb,amsthm}
\usepackage[T1]{fontenc}
\usepackage[margin=1.1in]{geometry}
\usepackage{booktabs}
\raggedbottom
\usepackage[hidelinks]{hyperref}

\newtheorem{theorem}{Theorem}[section]
\newtheorem{proposition}[theorem]{Proposition}
\newtheorem{lemma}[theorem]{Lemma}
\newtheorem{corollary}[theorem]{Corollary}
\theoremstyle{remark}
\newtheorem{remark}[theorem]{Remark}

\newcommand{\Z}{\mathbb Z}
\newcommand{\E}{\mathbb E}
\newcommand{\sleep}{\varrho}
\renewcommand{\Pr}{\mathbb P}

\title{Linear stabilization of subcritical activated random walk on the cycle}
\author{Yitzchak Shmalo}
\address{Einstein Institute of Mathematics, The Hebrew University of Jerusalem,\newline Givat Ram, Jerusalem, Israel}
\subjclass[2020]{Primary 60K35; Secondary 82C22}
\keywords{activated random walk, absorbing-state phase transition, stabilization time, odometer}
\date{}

\begin{document}

\begin{abstract}
Consider activated random walk on the cycle $\Z_n$ with symmetric nearest-neighbour jumps and sleep rate $\lambda>0$,
started from independent Poisson or Bernoulli particle counts of mean $\zeta$, all particles active. We show that if
$\zeta$ is below the critical density $\rho_c(\lambda)$, then the total number of topplings needed to stabilize the
system is of order $n$ with probability tending to one. Previously a bound of order $n\log^2 n$ was known in this range for a
fixed number of particles, from the work of Hoffman, Johnson and Junge. As a consequence the
odometer at a fixed site is tight. The proof combines the stable odometers built by
Johnson and Richey for the subcritical odometer tail on $\Z$ with a pasting argument on the cycle and a routing of
Poisson particles onto Bernoulli sites.
\end{abstract}

\maketitle

\section{Introduction}\label{sec:intro}

Activated random walk (ARW) is an interacting particle system with two kinds of particles, active and sleeping. Active
particles perform continuous-time random walks and fall asleep at rate $\lambda$ when alone on a site; a sleeping
particle is woken when another particle arrives at its site. On a finite graph without a sink, a configuration either
reaches a stable state, in which every particle sleeps, or keeps moving forever. We refer to the survey of Rolla
\cite{Rolla} for background.

We consider the model on the cycle $\Z_n=\Z/n\Z$ with nearest-neighbour symmetric jumps. Let $\eta_0$ consist of
independent Poisson$(\zeta)$ numbers of active particles at the sites of $\Z_n$, and let $T_n$ be the total number of
topplings performed during the stabilization of $\eta_0$, set to $+\infty$ if the system does not stabilize. Basu,
Ganguly, Hoffman and Richey \cite{BGHR} found two phases for this quantity, starting from independent Bernoulli$(\zeta)$
counts. They showed that $T_n\le Cn\log^2 n$ with probability at least $1-n^{-b}$ when $\zeta<\lambda/(1+\lambda)$
\cite[Theorem~1]{BGHR}, and that for every $\zeta\in(0,1)$, once $\lambda$ is small enough, $T_n\ge e^{cn}$ with
probability at least $1-e^{-cn}$ \cite[Theorem~2]{BGHR}. Rolla's survey \cite[Theorem~6.1]{Rolla} states both phases for
Poisson$(\zeta)$ counts, the fast one in the form $\Pr(T_n\le\kappa n\log^2 n)\ge 1-n^{-\delta}$ for large $\lambda$, with
a proof adapted from \cite{BGHR}. Among the open problems listed there are to show that one of the two phases occurs at
every point off the critical curve, to remove the factor $\log^2 n$ from the fast-phase bound, and to show that the
odometer at a fixed site is tight in the fast phase.

The first of these problems was settled by Hoffman, Johnson and Junge \cite{HJJ} as part of their proof of the density
conjecture; see also Forien \cite{Forien} for another proof of the density conjecture on the line. They identified the
critical density $\rho_c(\lambda)$ of ARW on $\Z$ with the growth rate of layer percolation, and showed
\cite[Proposition~8.12]{HJJ} that for $\lfloor\rho n\rfloor$ active particles placed uniformly on the cycle, the total
number of instructions used is at most $Cn\log^2 n$ with probability at least $1-n^{-b}$ when $\rho<\rho_c$, and at least
$e^{cn}$ with probability at least $1-e^{-cn}$ when $\rho>\rho_c$. Their statement is for a fixed number of particles;
when these are placed independently, its slow phase transfers to Poisson counts by the coupling used in the proof of
Corollary~\ref{cor:fixed} below, since the Poisson configuration then lies above a fixed number of independent particles
at a slightly smaller density. On the
slow side, Forien and Gaudilli\`ere \cite{FG} showed that for high density and small sleep rate the stabilization time on
the $d$-dimensional torus is exponentially large, in every dimension $d\ge1$, and de~Lima, Rolla and Terra \cite{LRT}
gave a toppling procedure for the ring at high density. Driven-dissipative ARW has also been studied: on a Euclidean ball its mixing time is at most
$1+o(1)$ times the volume as the radius tends to infinity \cite{LL}, on finite networks with
absorption at the boundary its relaxation time is known and separation cutoff is equivalent to the product condition
\cite{BS}, and in one dimension its mixing time has cutoff \cite{HJJMcut} and it shows self-organized criticality
\cite{HJJMsoc}. ARW on the complete graph is studied in \cite{JKM}. We show that the logarithmic factor can be removed in the
whole subcritical range. Recall that $\rho_c(\lambda)<1$ for every $\lambda>0$; this was proved on $\Z$ by Hoffman,
Richey and Rolla \cite{HRR} and in every dimension by Asselah, Forien and Gaudilli\`ere \cite{AFG}, so Bernoulli$(\zeta)$ counts
make sense below the critical density.

\begin{theorem}\label{thm:main}
Let $\lambda>0$ and $0<\zeta<\rho_c(\lambda)$. Let the initial counts on $\Z_n$ be independent with law either
Poisson$(\zeta)$ or Bernoulli$(\zeta)$, all particles active. There are constants $c,C,K\in(0,\infty)$, depending only
on $\lambda$ and $\zeta$, such that for all $n\ge 3$,
\begin{equation}\label{eq:main}
\Pr\bigl(cn\le T_n\le Kn\bigr)\ \ge\ 1-C\,\frac{\log n}{n}.
\end{equation}
\end{theorem}

The lower bound is elementary: each initially occupied site must topple at least once. For a similar
reason \cite{BGHR} note that their bound is sharp up to the polylogarithmic factor. The content of the theorem is the upper bound. The same bound holds for the fixed-number model of \cite[Proposition~8.12]{HJJ}.

\begin{corollary}\label{cor:fixed}
Let $\lambda>0$ and $0<\rho<\rho_c(\lambda)$, and place $\lfloor\rho n\rfloor$ active particles on $\Z_n$, either
independently and uniformly, or one at each site of a uniformly chosen set of $\lfloor\rho n\rfloor$ sites. Then \eqref{eq:main} holds for the total number of topplings $T_n$, with constants depending only
on $\lambda$ and $\rho$.
\end{corollary}

Write $m_n(x)$ for the number of topplings at $x$ during the stabilization of $\eta_0$, so that $T_n=\sum_x m_n(x)$.

\begin{corollary}\label{cor:tight}
Under the assumptions of Theorem~\ref{thm:main} or Corollary~\ref{cor:fixed},
\[
\lim_{M\to\infty}\limsup_{n\to\infty}\Pr\bigl(m_n(0)>M\bigr)=0 .
\]
\end{corollary}

Corollary~\ref{cor:tight} gives the tightness asked for in \cite[Section~6]{Rolla}, for the Poisson initial condition
considered there, at every density $\zeta<\rho_c(\lambda)$. When the particles in \cite[Proposition~8.12]{HJJ} are
placed independently, its slow phase transfers to Poisson counts by the coupling in the proof of Corollary~\ref{cor:fixed}, so that, apart from the critical density
itself, these are all the densities of the fast phase.

The main input is a construction of Johnson and Richey \cite{JR}. To prove that the odometer at the origin of
subcritical ARW on $\Z$ has a stretched-exponential tail, they build, with overwhelming probability, an odometer on an
interval that is stable in its interior, takes a value of order $t^2$ at the centre, and grows quadratically in the
distance from the centre. We use the growth at the two endpoints: it lets us paste local odometers centred at every site
of the cycle into a single function that dominates every legal toppling sequence (Section~\ref{sec:paste}). This gives
the upper bound for Bernoulli initial counts. Poisson counts are reduced to Bernoulli counts by first moving each
particle to a distinct site of an independent Bernoulli set of marks, at a cost of order $n$ topplings
(Section~\ref{sec:poisson}). The author proved in \cite{Shm} that for subcritical ARW on $\Z$
with i.i.d.\ initial counts the expected odometer at the origin is finite exactly when the count per site has a finite
third moment; the upper bound there uses a similar routing onto an independent Bernoulli configuration of slightly
larger density. That result concerns the odometer on $\Z$ and gives no bound on the stabilization time of the cycle.

\section{Notation and the site-wise representation}\label{sec:prelim}

Throughout, $c,C,C',\ldots$ denote positive constants that depend only on $\lambda$ and the particle density; their
values may change from line to line.

We use the site-wise representation of ARW \cite[Section~2]{Rolla}. A configuration is a function
$\eta:\Z_n\to\{0,\sleep,1,2,\ldots\}$, where $\eta(x)=\sleep$ means one sleeping particle at $x$ and $\eta(x)=k\ge1$
means $k$ active particles. We order $0<\sleep<1<2<\cdots$, and write $|\sleep|=1$ for the number of particles. A
configuration is stable if every value is $0$ or $\sleep$.

Each site $x$ carries an instruction stack $(\tau^{x,k})_{k\ge1}$. The instructions are independent over $x$ and $k$;
each is a jump to the left neighbour or to the right neighbour with probability $\frac{1}{2(1+\lambda)}$ each, or a sleep
instruction with probability $\frac{\lambda}{1+\lambda}$. Toppling $x$ applies the next unused instruction at $x$: a jump
moves one particle from $x$ to the neighbour (a sleeping particle at $x$ is woken first, and the moved particle is active
at its destination, waking a sleeping particle there); a sleep instruction turns $\eta(x)=1$ into $\sleep$ and does
nothing if $\eta(x)\ge2$ or $\eta(x)=\sleep$. A toppling at $x$ is \emph{legal} if $\eta(x)\ge1$ and \emph{acceptable} if
$\eta(x)\ge\sleep$. For a finite sequence $\alpha$ of topplings, $m_\alpha(x)$ denotes the number of topplings of $x$ in
$\alpha$.

We use three properties from \cite[Section~2]{Rolla}: the abelian property (Lemma~2.4 there), the least-action principle
(Lemma~2.1: if $\alpha$ is an acceptable sequence that stabilizes $\eta$ and $\beta$ is any legal sequence for $\eta$,
then $m_\beta\le m_\alpha$), and monotonicity (Lemma~2.5: if $\eta\le\tilde\eta$ pointwise then
$m_\eta\le m_{\tilde\eta}$, where $m_\eta$ is the stabilizing odometer). They are stated in \cite{Rolla} for a finite set
$V\subset\Z^d$ with particles leaving $V$ frozen outside, and their proofs use only the finiteness of the set of sites
and the site-wise structure; they apply verbatim to the finite graph $\Z_n$ with $V=\Z_n$. In particular, if every legal
sequence for $\eta$ has length at most $L$, then a maximal legal sequence stabilizes $\eta$ and $T_n=\sum_x m_\eta(x)\le L$.

\subsection*{Stable arrays}
For a function $u:I\to\Z_{\ge0}$ on an interval $I$ of sites and $y\in I$, let $L_y(u)$ and $R_y(u)$ be the numbers of
left and right jump instructions among $\tau^{y,1},\ldots,\tau^{y,u(y)}$, and let $s_y(u)=1$ if $u(y)\ge1$ and
$\tau^{y,u(y)}$ is a sleep instruction, and $s_y(u)=0$ otherwise. For an initial configuration $\sigma$ of active
particles, we say that $u$ is \emph{stable at} $y$ if $y-1,y+1\in I$ and
\begin{equation}\label{eq:balance}
\sigma(y)+R_{y-1}(u)+L_{y+1}(u)-L_y(u)-R_y(u)=s_y(u).
\end{equation}
This is the mass-balance equation (2.1) of \cite{JR}: the left side is the number of particles that $u$ leaves at $y$,
and stability means that this number is $0$, or $1$ with the last instruction executed at $y$ a sleep instruction. We
call $u$ stable in the interior of $I=[x-R,x+R]$ if it is stable at every $y$ with $|y-x|<R$. The values of $u$ at the
two endpoints are unconstrained.

\section{Local odometers with large endpoints}\label{sec:local}

\begin{lemma}\label{lem:local}
Let $\lambda>0$ and $0<r<\rho_c(\lambda)$, and let $\sigma$ consist of independent Bernoulli$(r)$ active particles.
There are constants $a,b,c,C>0$ depending only on $\lambda$ and $r$ such that for all integers $R>t\ge1$, with
probability at least $1-Ce^{-ct}$ there is a function $u:[-R,R]\to\Z_{\ge0}$, stable in the interior of $[-R,R]$, with
\begin{equation}\label{eq:local}
u(0)\le bt^2\qquad\text{and}\qquad \min\{u(-R),u(R)\}\ge aR^2 .
\end{equation}
The event depends only on $\sigma$ and the instruction stacks on $[-R,R]$.
\end{lemma}

\begin{proof}
This is the construction in the proof of the upper bound of \cite[Theorem~1.1]{JR}, applied with their $n$
equal to $t$, their $N$ equal to $R$, their $\rho$ equal to $r$, and $\epsilon=(\rho_*-r)/3$, where $\rho_*=\rho_c(\lambda)$
by \cite{HJJ}. The lemma is stated for Bernoulli counts because that construction is carried out for i.i.d.
Bernoulli initial configurations; in particular the net flow $f_0=\sigma(0)$ below lies in $\{0,1\}$. We recall what it
gives. Put $u_0=\lceil 3(1+\lambda)t^2\rceil$ and $f_0=\sigma(0)$. Johnson and Richey take
the extended odometer $g$ on $[0,R]$ with $g(0)=u_0$, net flow $f_0$ from $0$ to $1$, stable on $[1,R-1]$, that
corresponds to the $k_0$-greedy infection path of layer percolation. By \cite[Lemma~3.2(a)]{JR} and the minimality of the
minimal odometer, $g\ge0$ on $[0,t]$ except with probability $Ce^{-ct}$. For $t<j\le R$ the right-jump count of $g$ at $j$
equals $r_j+R_j(m)$, where $r_j$ is the first coordinate of the greedy path and $m$ the minimal odometer
\cite[Proposition~2.1]{JR}, and \cite[Lemmas~3.2(b) and~3.3]{JR} give
\[
\Pr\Bigl(R_j(g)<(\rho_*-r-2\epsilon)\,j^2/2\Bigr)\le Ce^{-cj}.
\]
Negative values of an extended odometer correspond to executing instructions in reverse \cite[Section~2]{JR}, following the definition in \cite{HJJ}, so a
positive right-jump count at $j$ forces $g(j)>0$. A union bound over $t<j\le R$ shows that, except with probability $Ce^{-ct}$, $g$ is a nonnegative odometer on $[0,R]$
stable on $[1,R-1]$ with
\[
g(R)\ \ge\ R_R(g)\ \ge\ \tfrac12(\rho_*-r-2\epsilon)R^2=\tfrac{\epsilon}{2}R^2 .
\]
By the symmetry between left and right instructions, the same holds for an odometer $g'$ on $[-R,0]$ with $g'(0)=u_0$,
zero net flow from $0$ to $-1$, stable on $[-R+1,-1]$, and $g'(-R)\ge\frac{\epsilon}{2}R^2$. Pasting $g$ and $g'$ along
the common value $u_0$ at $0$ and lowering the value at $0$ to the last non-sleep instruction among the first $u_0$
(which exists except with probability $Ce^{-ct}$), Johnson and Richey obtain an odometer on $[-R,R]$ that is stable on
$[-R+1,R-1]$ (same proof in \cite{JR}). Lowering the value at $0$ changes neither the jump counts at $0$ nor the values elsewhere.
The resulting function satisfies \eqref{eq:local} with $b=3(1+\lambda)+1$ and $a=\epsilon/2$. The constants in these
bounds do not depend on $R$, which plays the role of their $N$.

Finally, the existence of a function with the stated properties is an event determined by $\sigma$ and the stacks on
$[-R,R]$, since it is a countable union over candidate functions, each of which is checked by finitely many
instructions.
\end{proof}

\section{Pasting on the cycle: the Bernoulli case}\label{sec:paste}

The following comparison is a version of the least-action principle for functions satisfying the mass-balance relation
as an inequality; compare the least-action principle for weakly stable odometers in \cite[Lemma~2.3]{HJJ}.

\begin{lemma}\label{lem:super}
Let $n\ge3$, let $\sigma$ be a configuration of active particles on $\Z_n$, and let $U:\Z_n\to\Z_{\ge0}$ satisfy
\begin{equation}\label{eq:super}
\sigma(y)+R_{y-1}(U)+L_{y+1}(U)-L_y(U)-R_y(U)\ \le\ s_y(U)\qquad\text{for every }y\in\Z_n .
\end{equation}
Then every legal sequence $\beta$ for $\sigma$ satisfies $m_\beta\le U$. Consequently $T_n\le\sum_y U(y)$.
\end{lemma}

\begin{proof}
Suppose not, and consider the first toppling of $\beta$ after which $m_\beta(y)>U(y)$ at some site $y$. Just before it,
$y$ has been toppled exactly $U(y)$ times and every other site $z$ at most $U(z)$ times. The number of particles at $y$ at
that moment is $\sigma(y)$ plus the jumps into $y$ minus the jumps out of $y$, that is
\[
\sigma(y)+R_{y-1}(m)+L_{y+1}(m)-L_y(U)-R_y(U),
\]
where $m$ is the current odometer. Since $R_{y-1}$ and $L_{y+1}$ are nondecreasing in the number of instructions used,
this is at most the left side of \eqref{eq:super}, hence at most $s_y(U)\in\{0,1\}$. If $s_y(U)=0$, the site $y$ is empty
and cannot be toppled legally. If $s_y(U)=1$, the last instruction used at $y$ was a sleep instruction, used by a
legal toppling, so $y$ held at least one particle at that time. If at that time
$y$ held at least two particles, it still holds at least two, since no instruction has been used at $y$ since then and
arrivals only add particles; this contradicts the bound $1$. If it held exactly one particle, that particle fell asleep,
and it is still alone and asleep, since an arrival would raise the count to two. In both cases the next toppling at $y$
is not legal, a contradiction. The last statement follows from the remark at the end of Section~\ref{sec:prelim}.
\end{proof}

\begin{proposition}\label{prop:bernoulli}
Theorem~\ref{thm:main} holds for Bernoulli$(r)$ initial counts with $0<r<\rho_c(\lambda)$.
\end{proposition}

\begin{proof}
Let $a,b,c,C$ be the constants of Lemma~\ref{lem:local}. Set $J=\lceil B\log n\rceil$ and $R=\lceil AJ\rceil$, where $A>1$
satisfies $aA^2>b$ and $B$ is large, to be fixed. We take $n>2R+1$, so that every interval $[x-R,x+R]$ of $\Z_n$ is a
genuine interval; smaller $n$ are absorbed into $C$.

For each $x\in\Z_n$ let $D_x$ be the least $t\in\{1,\ldots,J\}$ for which there is a function on $[x-R,x+R]$ satisfying
the conclusion of Lemma~\ref{lem:local} translated to centre $x$, and $D_x=J$ if there is none. Let $G$ be the event that
such a function exists at every centre for some $t\le J$. By Lemma~\ref{lem:local} and translation invariance,
$\Pr(D_x>t)\le Ce^{-ct}$ for $t<J$, and $\Pr(G^c)\le nCe^{-cJ}\le n^{-2}$ once $B$ is large. In particular
$\sup_{n,x}\E D_x^4<\infty$, and $D_x$ is determined by $\sigma$ and the stacks within distance $R$ of $x$.

On $G$, choose for each centre $x$ a function $u_x$ on $[x-R,x+R]$ with $u_x(x)\le bD_x^2$, stable in the interior,
with endpoint values at least $aR^2$. For $y\in\Z_n$ put
\[
U(y)=\min\{u_x(y):\ |x-y|\le R\},
\]
with cyclic distance. We check \eqref{eq:super} at $y$. Let $x$ attain the minimum. Then
\[
u_x(y)=U(y)\le u_y(y)\le bD_y^2\le bJ^2<aR^2\le\min\{u_x(x-R),u_x(x+R)\},
\]
so $y$ is not an endpoint of $[x-R,x+R]$, and $u_x$ is stable at $y$. The neighbours $y\pm1$ lie in $[x-R,x+R]$, so
$U(y\pm1)\le u_x(y\pm1)$. Since $U(y)=u_x(y)$, the quantities $L_y,R_y,s_y$ agree for $U$ and $u_x$, while $R_{y-1}(U)\le
R_{y-1}(u_x)$ and $L_{y+1}(U)\le L_{y+1}(u_x)$. Hence \eqref{eq:balance} for $u_x$ gives \eqref{eq:super} for $U$ at $y$.
By Lemma~\ref{lem:super}, on $G$,
\[
T_n\le\sum_yU(y)\le b\sum_y D_y^2 .
\]
If the cyclic distance between $x$ and $x'$ exceeds $2R$, then $D_x$ and $D_{x'}$ are independent. By the Cauchy--Schwarz
inequality, $\operatorname{Var}(\sum_yD_y^2)\le(4R+1)\,n\sup\E D^4\le C'n\log n$, while $\E\sum_yD_y^2\le C_0n$. By
Chebyshev's inequality, $\Pr(\sum_yD_y^2>(C_0+1)n)\le C'\log n/n$. Therefore $T_n\le b(C_0+1)n$ except with probability
$C''\log n/n$. The lower bound is proved in Section~\ref{sec:lower}.
\end{proof}

\section{Poisson counts}\label{sec:poisson}

Fix $r$ with $\zeta<r<\rho_c(\lambda)$, and let $\xi$ be an independent set of marks: each site of $\Z_n$ is marked with
probability $r$, independently. We first move every particle of $\eta_0$ to a distinct mark by acceptable topplings, and
then stabilize.

\subsection*{The routing}
Enumerate the particles of $\eta_0$ in a fixed order. In turn, each particle is moved until it stands on a mark that has
not been reserved, which it then reserves for good; its starting site counts. To move a particle standing at $z$ we
topple $z$. A jump instruction moves this particle (any particle at $z$ may be chosen, since particles are
indistinguishable in the site-wise representation, and we follow the one being routed); a sleep instruction leaves it in
place, possibly asleep, and we topple $z$ again, so that the first later jump instruction at $z$ moves it. Every such
toppling is acceptable, since $z$ holds the routed particle. Particles already
standing on their reserved marks are not followed: they stay at their sites, possibly asleep, possibly woken by a passing
particle.

It is convenient to generate the instructions used during the routing from tapes attached to the particles. Index the
particles by their starting site and their rank among the particles starting there: particle $(x,k)$ carries a tape
$\theta^{x,k}$, a sequence of instructions with the same law as a stack, the tapes being independent over $(x,k)$ and
independent of $\eta_0$ and $\xi$. Each time the site of a particle is toppled for it, the next symbol of its own tape is used as the next
instruction of the site where it stands. Since each symbol is independent of
everything revealed before it is used, the sequence of instructions used at each site has the law of the beginning of
an instruction stack, and the unused instructions at each site may be taken to be fresh independent stacks, independent of
$\eta_0$, $\xi$ and all tapes. We call these the \emph{fresh stacks}.

Under the tape description, each routed particle performs a lazy simple random walk driven by its tape: it jumps left or right
with probability $\frac{1}{2(1+\lambda)}$ each and stays with probability $\frac{\lambda}{1+\lambda}$ at each step, until
it reaches an unreserved mark. On the cycle such a walk reaches any given site almost surely, so the routing terminates
whenever the number of marks exceeds the number of particles.

\subsection*{The cost of the routing}
Call a set of consecutive sites of $\Z_n$ \emph{deficient} if it contains at least as many particles of $\eta_0$ as
marks; the whole cycle is included among these sets. For a set of $l$ sites, the difference between the particle count
and the mark count is a sum of $l$ independent copies of $X-Y$, with $X\sim$ Poisson$(\zeta)$ and $Y\sim$ Bernoulli$(r)$
independent, and $\E[X-Y]=\zeta-r<0$. By Chernoff's bound there is $d>0$ such that every fixed set of $l$ consecutive sites is
deficient with probability at most $e^{-dl}$. Let $H=\lceil B'\log n\rceil$, and let $\mathcal D$ be the event that no
deficient set has more than $H$ sites. Since there are at most $n$ sets of each length,
$\Pr(\mathcal D^c)\le n\sum_{l>H}e^{-dl}\le n^{-2}$ for $B'$ large. On $\mathcal D$ the whole cycle is not deficient,
so the routing terminates.

Assume $\mathcal D$. Let $\mathcal V$ be the set of sites visited by at least one routed particle, starting and final
sites included, and let $I$ be a connected component of $\mathcal V$; it is a set of consecutive sites. A routed particle
visits a set of consecutive sites containing its starting site, so a particle starting in $I$ stays in $I$, and a particle
that ends in $I$ started in $I$. Every mark in $I$ has been visited, and a particle visiting an unreserved mark stops
there, so every mark in $I$ is reserved, each by a different particle that started in $I$. Hence $I$ contains at least as
many particles of $\eta_0$ as marks: $I$ is deficient, and so has at most $H$ sites.

For $x\in\Z_n$ let $h_x$ be the largest number of sites of a deficient set with at most $H$ sites containing $x$, or $1$
if there is none. A deficient set of $l$ sites containing $x$ is one of $l$ sets, so $\Pr(h_x\ge k)\le\sum_{l\ge k}le^{-dl}$
for $k\ge2$, uniformly in $n$ and $x$; all moments of $h_x$ are bounded uniformly. On $\mathcal D$, a particle starting at
$x$ moves inside its component, which contains $x$ and has at most $h_x$ sites, so its walk does not leave
$[x-h_x,x+h_x]$. Let $W_x$ be the sum, over the particles starting at $x$, of the number of steps their tape-driven walks
take to leave $[x-h_x,x+h_x]$, each walk continued on its own tape for this purpose. On $\mathcal D$ the number of
instructions used to route the particles starting at $x$ is at most $W_x$.

For the lazy walk above started at $i\in[-h,h]$, the mean exit time from $[-h,h]$ is $(1+\lambda)((h+1)^2-i^2)$, by
solving the difference equation with zero boundary values at $\pm(h+1)$. By Markov's inequality and restarting at
multiples of $\lceil2(1+\lambda)(h+1)^2\rceil$ steps, its second moment is at most $C_\lambda(h+1)^4$. The tapes are
independent of $\eta_0$ and $\xi$, so by the Cauchy--Schwarz inequality,
$\E[W_x^2\mid\eta_0,\xi]\le C_\lambda\,\eta_0(x)^2(h_x+1)^4$, and another application with the fourth moment of the
Poisson law and the eighth moment of $h_x+1$ gives $\sup_{n,x}\E W_x^2<\infty$. The variable $W_x$ is determined by $\eta_0$
and $\xi$ within distance $H$ of $x$ and by the tapes $\theta^{x,k}$, $k\ge1$, so $W_x$ and $W_{x'}$ are
independent when the cyclic distance between $x$ and $x'$ exceeds $2H$. As in the proof of
Proposition~\ref{prop:bernoulli}, Chebyshev's inequality gives $\sum_xW_x\le C_1n$ except with probability
$C'\log n/n$.

\subsection*{After the routing}
When the routing ends, each particle stands on its own mark, active or asleep, and there is at most one particle at
each site. The resulting configuration $\eta'$ therefore satisfies $\eta'\le\eta_\xi$, where $\eta_\xi$ has one active
particle at each mark. Apply Proposition~\ref{prop:bernoulli} to $\eta_\xi$ and the fresh stacks: since $\xi$ is
Bernoulli$(r)$ and independent of the fresh stacks, the stabilization of $\eta_\xi$ uses at most $C_2n$ topplings except
with probability $C\log n/n$. By monotonicity, the same holds for the stabilization of $\eta'$ with the fresh stacks.

The routing followed by a legal stabilization of $\eta'$ with the fresh stacks is an acceptable sequence of topplings
that stabilizes $\eta_0$ with the original stacks. By the least-action principle, every legal sequence for $\eta_0$, in
particular a stabilizing one, uses at most as many topplings at each site. Collecting the exceptional events,
$T_n\le(C_1+C_2)n$ except with probability $C\log n/n$. This proves the upper bound in Theorem~\ref{thm:main} for Poisson
counts.

\section{The lower bound and the corollaries}\label{sec:lower}

\begin{proof}[Proof of the lower bound in Theorem~\ref{thm:main}]
An initially occupied site holds only active particles, and the configuration can become stable at that site only after
an instruction has been used there. Hence $T_n$ is at least the number of initially occupied sites, which is binomial with
parameters $n$ and $p=1-e^{-\zeta}$ or $p=\zeta$. By Chernoff's bound it is at least $pn/2$ except with probability
$e^{-c n}$.
\end{proof}

\begin{proof}[Proof of Corollary~\ref{cor:fixed}]
Let $N=\lfloor\rho n\rfloor$ and choose $\zeta\in(\rho,\rho_c(\lambda))$. Take independent Poisson$(\zeta)$ counts on
$\Z_n$, with total $M$. Given $M$, these counts are those of $M$ independent uniform particles, and when $M\ge N$ the
first $N$ of these particles give a configuration with the law in the corollary that is pointwise below the Poisson
configuration. Since $\Pr(M<N)\le e^{-cn}$, monotonicity and Theorem~\ref{thm:main} give the upper bound. For the lower
bound, the number of sites occupied by $N$ independent uniform particles has mean $n(1-(1-1/n)^N)\ge c'n$ and changes by at
most one when one particle is moved, so it is at least $c'n/2$ except with probability $e^{-cn}$ by the bounded
differences inequality.

For a uniformly chosen set of $N$ sites, take Bernoulli$(\zeta)$ counts instead, with the same $\zeta$. Given the number
$M'$ of occupied sites, the occupied set is uniform among sets of that size, so when $M'\ge N$ a uniformly chosen
$N$-element subset of it has the law in the corollary and is pointwise below the Bernoulli configuration. Since
$\Pr(M'<N)\le e^{-cn}$, monotonicity and Theorem~\ref{thm:main} for Bernoulli counts give the upper bound, and the lower
bound holds because $N$ sites are occupied.
\end{proof}

\begin{proof}[Proof of Corollary~\ref{cor:tight}]
The law of $(m_n(x))_x$ is invariant under rotations of $\Z_n$. On the event $T_n\le Kn$ at most $Kn/M$ sites have
$m_n(x)>M$. Hence
\[
\Pr(m_n(0)>M)=\frac1n\,\E\#\{x:m_n(x)>M\}\le\frac KM+\Pr(T_n>Kn),
\]
and the last term tends to zero by Theorem~\ref{thm:main} or Corollary~\ref{cor:fixed}.
\end{proof}

\begin{remark}
Theorem~\ref{thm:main} is a statement with high probability, not about expectations: for Poisson counts there is a
positive probability that the cycle holds more than $n$ particles, in which case $T_n=\infty$. The rate $\log n/n$ comes
from the second-moment bounds and is not optimized.
\end{remark}

\section*{Acknowledgments and funding}

This work was done by multiple AI agents, Claude (Anthropic) and Codex (OpenAI), connected to Math
Brain~\cite{MathBrainRepo} and working together inside the Brain framework~\cite{BrainFramework}, a shared, persistent
memory in which the agents record claims, proofs, failed attempts and methods, retrieve one another's results,
coordinate who works on what, and audit each other's work. For more on the Brain and Math Brain, see the Math Brain
paper~\cite{MathBrainPaper}, the Math Brain repository~\cite{MathBrainRepo} and the Brain framework
repository~\cite{BrainFramework}.

The research presented in this paper was supported by the European Research Council (ERC) under the European
Union's Horizon Europe research and innovation programme (grant agreement No.\ 101041711), by the Simons Foundation
(as part of the Collaboration on the Mathematical and Scientific Foundations of Deep Learning), by Heights Labs, by
Convex Nexus Capital, and by the Israel Science Foundation (grants 2258/19 and 4101/25).

\section*{Code and data availability}

The paper uses no data and no code.

\section*{Declaration of competing interest}

The author holds equity in Heights Labs and Convex Nexus Capital and has a management role at Convex Nexus Capital.

\section*{Declaration of generative AI and AI-assisted technologies}

This paper is part of a broader research project on proving open problems with AI systems. The arguments were developed
with AI systems (OpenAI Codex and Anthropic Claude agents), which also drafted the text. The problem is posed in Rolla's
survey \cite{Rolla}, and the two phases on the cycle are due to Basu, Ganguly, Hoffman and Richey \cite{BGHR}. The proof rests on the stable odometers constructed by Johnson and Richey \cite{JR} and on the
identification of the critical density by Hoffman, Johnson and Junge \cite{HJJ}. The author directed the work and is
responsible for the content of the paper and for decisions concerning its dissemination.

\appendix
\section{Notation}

\begin{center}
\begin{tabular}{@{}lp{0.52\textwidth}l@{}}
\toprule
Symbol & Meaning & First appears\\
\midrule
$\Z_n$ & the cycle with $n$ sites & Section~\ref{sec:intro}\\
$\lambda$ & sleep rate & Section~\ref{sec:intro}\\
$\zeta$ & mean initial count per site & Section~\ref{sec:intro}\\
$\rho_c(\lambda)$ & critical density of ARW on $\Z$ & Section~\ref{sec:intro}\\
$T_n$ & total number of topplings during stabilization & Section~\ref{sec:intro}\\
$c,C,K$ & positive constants depending only on $\lambda$ and the density & Theorem~\ref{thm:main}\\
$m_n(x)$ & number of topplings at $x$ & Corollary~\ref{cor:tight}\\
$\sleep$ & one sleeping particle & Section~\ref{sec:prelim}\\
$\tau^{x,k}$ & $k$-th instruction at site $x$ & Section~\ref{sec:prelim}\\
$m_\alpha$, $m_\eta$ & odometer of a sequence of topplings; stabilizing odometer of $\eta$ & Section~\ref{sec:prelim}\\
$L_y(u)$, $R_y(u)$ & left and right jumps among the first $u(y)$ instructions at $y$ & Section~\ref{sec:prelim}\\
$s_y(u)$ & $1$ if the $u(y)$-th instruction at $y$ is a sleep instruction, else $0$ & Section~\ref{sec:prelim}\\
$\sigma$ & initial configuration of active particles & Section~\ref{sec:prelim}\\
$a,b$ & endpoint and centre constants of the local odometers & Lemma~\ref{lem:local}\\
$D_x$, $J$, $R$ & local scale at $x$, its cap, the window radius & Proposition~\ref{prop:bernoulli}\\
$U$ & pasted odometer & Proposition~\ref{prop:bernoulli}\\
$\xi$, $r$ & Bernoulli marks and their density & Section~\ref{sec:poisson}\\
$H$, $h_x$ & cap on the size of deficient sets; largest deficient set at $x$ & Section~\ref{sec:poisson}\\
$\theta^{x,k}$ & tape of the $k$-th particle starting at $x$ & Section~\ref{sec:poisson}\\
$W_x$ & routing cost of the particles starting at $x$ & Section~\ref{sec:poisson}\\
\bottomrule
\end{tabular}
\end{center}

\begin{thebibliography}{99}\raggedright

\bibitem{AFG}
A.~Asselah, N.~Forien and A.~Gaudilli\`ere,
\emph{The critical density for activated random walks is always less than 1},
Ann. Probab. \textbf{52} (2024), no.~5, 1607--1649, doi:10.1214/23-AOP1674.

\bibitem{BGHR}
R.~Basu, S.~Ganguly, C.~Hoffman and J.~Richey,
\emph{Activated random walk on a cycle},
Ann. Inst. Henri Poincar\'e Probab. Stat. \textbf{55} (2019), no.~3, 1258--1277, doi:10.1214/18-AIHP918;
arXiv:1709.09163.

\bibitem{BS}
A.~Bristiel and J.~Salez,
\emph{Separation cutoff for activated random walks},
Ann. Appl. Probab. \textbf{34} (2024), no.~6, 5211--5227, doi:10.1214/24-AAP2089; arXiv:2209.03274.

\bibitem{LRT}
B.~N.~B. de~Lima, L.~T. Rolla and C.~Terra,
\emph{On the slow phase for fixed-energy activated random walks},
J. Appl. Probab. \textbf{63} (2026), no.~3, 1330--1345, doi:10.1017/jpr.2026.10083; arXiv:2408.07049.

\bibitem{Forien}
N.~Forien,
\emph{A new proof of superadditivity and of the density conjecture for activated random walks on the line},
preprint (2025), arXiv:2502.02579.

\bibitem{FG}
N.~Forien and A.~Gaudilli\`ere,
\emph{Active phase for activated random walks on the lattice in all dimensions},
Ann. Inst. Henri Poincar\'e Probab. Stat. \textbf{60} (2024), no.~2, 1188--1214, doi:10.1214/22-AIHP1341;
arXiv:2203.02476.

\bibitem{HJJ}
C.~Hoffman, T.~Johnson and M.~Junge,
\emph{The density conjecture for activated random walk},
preprint (2024), arXiv:2406.01731.

\bibitem{HJJMcut}
C.~Hoffman, T.~Johnson, M.~Junge and J.~Meisel,
\emph{Cutoff for activated random walk},
preprint (2025), arXiv:2501.17938.

\bibitem{HJJMsoc}
C.~Hoffman, T.~Johnson, M.~Junge and J.~Meisel,
\emph{Activated random walk exhibits self-organized criticality},
preprint (2026), arXiv:2605.00207.

\bibitem{HRR}
C.~Hoffman, J.~Richey and L.~T. Rolla,
\emph{Active phase for activated random walk on $\Z$},
Comm. Math. Phys. \textbf{399} (2023), no.~2, 717--735, doi:10.1007/s00220-022-04572-x; arXiv:2009.09491.

\bibitem{JR}
T.~Johnson and J.~Richey,
\emph{The odometer in subcritical activated random walk},
Electron. Commun. Probab. \textbf{31} (2026), paper no.~74, doi:10.1214/26-ECP819; arXiv:2510.05514.

\bibitem{JKM}
M.~Junge, H.~Kaufman and J.~Meisel,
\emph{Scaling limit and density conjecture for activated random walk on the complete graph},
preprint (2026), arXiv:2604.04747.

\bibitem{LL}
L.~Levine and F.~Liang,
\emph{Exact sampling and fast mixing of activated random walk},
Electron. J. Probab. \textbf{29} (2024), 1--20, doi:10.1214/24-EJP1220; arXiv:2110.14008.

\bibitem{Rolla}
L.~T. Rolla,
\emph{Activated random walks on $\mathbb Z^d$},
Probab. Surv. \textbf{17} (2020), 478--544, doi:10.1214/19-PS339; arXiv:1906.05037.

\bibitem{Shm}
Y.~Shmalo,
\emph{A third-moment criterion for the odometer of subcritical activated random walk on $\mathbb Z$},
Zenodo (2026), doi:10.5281/zenodo.24363818.

\bibitem{MathBrainPaper}
Y.~Shmalo, \emph{Math Brain: shared, auditable memory for multi-agent mathematical research}, October 2026.
\url{https://github.com/yspennstate/math-brain-public/tree/5427734deee9be643d6dfbdd5e7d4baeeefe4a4a/papers/math-brain-2026-10-24}

\bibitem{MathBrainRepo}
Y.~Shmalo, \emph{Math Brain: public snapshot, code and tests}, 2026.
\url{https://github.com/yspennstate/math-brain-public/tree/5427734deee9be643d6dfbdd5e7d4baeeefe4a4a}

\bibitem{BrainFramework}
Y.~Shmalo, \emph{Brain Framework}, public software repository, 2026.
\url{https://github.com/yspennstate/brain-framework-public/tree/46b935d26d5a2893f7790b60466387dd8da47ef5}

\end{thebibliography}

\end{document}
